\subsection{General Properties of Quaternion Algebras}

Let \(\alpha, \beta, \gamma\) be elements of K, and let F be the quaternion algebra of type \((\alpha, \beta, \gamma)\) . Recall (III, §2, No.~5, p.~445) \(^{(1)}\) that F is an associative unital K-algebra that has a basis \((1, i, j, k)\) over K satisfying the relations

\[
i ^ {2} = \alpha + \beta i, j ^ {2} = \gamma , i j = k, j i = \beta j - k.\tag{1}
\]

It is a Cayley algebra (III, §2, No.~4, p.~441, Definition 1) whose conjugation satisfies

\[
\overline {{i}} = \beta - i, \quad \overline {{j}} = - j, \quad \overline {{k}} = - k.\tag{2}
\]

Recall that the Cayley trace and norm of F are the mappings \(T_{F}\) and \(N_{F}\) from F to K defined by \(\mathrm{T}_{\mathrm{F}}(q)=q+\overline{q}\) and \(\mathrm{N}_{\mathrm{F}}(q)=q\overline{q}\) .

The linear subspace E of F with basis \((1,i)\) is a commutative Cayley subalgebra of F; it is a quadratic algebra of type \((\alpha,\beta)\) , and F can be identified with the Cayley extension of E defined by \(\gamma\) (III, §2, No.~5, p.~444). For every \(z\in E\) , we have \(zj=j\overline{z}\) . Every element q of F can be written uniquely as \(x+jy\) with \(x,y\in E\) , and we have

\[
\overline {{q}} = \overline {{x}} - j y, \quad \mathrm{T} _ {\mathrm{F}} (q) = x + \overline {{x}}, \quad \mathrm{N} _ {\mathrm{F}} (q) = x \overline {{x}} - \gamma y \overline {{y}}.\tag{3}
\]

PROPOSITION 1. --- The characteristic polynomial of an element q of F is equal to \(\left(\mathrm{X}^{2}-\mathrm{T}_{\mathrm{F}}(q)\mathrm{X}+\mathrm{N}_{\mathrm{F}}(q)\right)^{2}\) .

By the above, the algebra \(\mathrm{F}\) is a free right \(\mathrm{E}\) -module with basis \((1,j)\) . Consequently, \(\mathrm{F}[\mathrm{X}]\) is a free right \(\mathrm{E}[\mathrm{X}]\) -module with basis \((1,j)\) . We denote by \(u\) the endomorphism of the right \(\mathrm{E}[\mathrm{X}]\) -module \(\mathrm{F}[\mathrm{X}]\) defined by \(u(\mathrm{P}) = (\mathrm{X} - q)\mathrm{P}\) for \(\mathrm{P} \in \mathrm{F}[\mathrm{X}]\) . The characteristic polynomial of \(q\) is the determinant of \(u\) viewed as an endomorphism of the \(\mathrm{K}[\mathrm{X}]\) -module \(\mathrm{F}[\mathrm{X}]\) . By Proposition 6 of III, §9, No.~4, p.~546, it is equal to \(\mathrm{N}(\det u)\) , where \(\mathrm{N}\) denotes the norm from \(\mathrm{E}[\mathrm{X}]\) to \(\mathrm{K}[\mathrm{X}]\) . Let us write \(q\) as \(x + jy\) with \(x, y \in \mathrm{E}\) . The matrix of \(u\) with respect to the basis \((1,j)\) is \(\begin{pmatrix} \mathrm{X} - x & -\gamma \overline{y} \\ -y & \mathrm{X} - \overline{x} \end{pmatrix}\) ; its determinant is equal to \(\mathrm{D} = (\mathrm{X} - x)(\mathrm{X} - \overline{x}) - \gamma y\overline{y} = \mathrm{X}^2 - \mathrm{T}_{\mathrm{F}}(q)\mathrm{X} + \mathrm{N}_{\mathrm{F}}(q)\) (cf.~formula (3)). Since \(\mathrm{D}\) belongs to \(\mathrm{K}[\mathrm{X}]\) , we have \(\mathrm{N}(\mathrm{D}) = \mathrm{D}^2\) ; Proposition 1 follows.

Remarks. --- 1) Suppose that the characteristic of K is different from 2, and set \(i' = 2i - \beta\) . Then \((1, i')\) is a basis of E over K, and we have \(i'^2 = 4\alpha + \beta^2\) . It follows that E is isomorphic to the quadratic algebra of type \((4\alpha + \beta^2, 0)\) and F to the quaternion algebra of type \((4\alpha + \beta^2, 0, \gamma)\) .

\begin{enumerate}
\def\labelenumi{\arabic{enumi})}
\setcounter{enumi}{1}
\item
  The quaternion algebra of type \((\alpha,0,\gamma)\) is isomorphic to the quaternion algebra of type \((\gamma,0,\alpha)\) (III, §2, No.~5, p.~445). It is also isomorphic to the quaternion algebra of type \((\alpha a^{2},0,\gamma c^{2})\) for every pair \((a,c)\) of nonzero elements of K.
\item
  Let q be an element of F. Then q is nilpotent if and only if its characteristic polynomial is equal to \(X^{4}\) , that is, \(\mathrm{T}_{\mathrm{F}}(q)=\mathrm{N}_{\mathrm{F}}(q)=0\) ; we then have \(q^{2}=0\) .
\end{enumerate}

Example. --- The matrix algebra \(\mathbf{M}_{2}(\mathrm{K})\) is isomorphic to the quaternion algebra of type \((0,1,1)\) . Indeed, consider the quadratic algebra \(E = K \times K\) (of type \((0,1)\) ) and the quaternion algebra \(F = E + Ej\) , which is the Cayley extension of E defined by the element \(\gamma = 1\) . The mapping \((a,b) \mapsto \begin{pmatrix} a & 0 \\ 0 & b \end{pmatrix}\) is an algebra homomorphism from E to \(\mathbf{M}_{2}(\mathrm{K})\) . Since for a, b in K, we have

\[
\left( \begin{array}{c c} 0 & 1 \\ 1 & 0 \end{array} \right) \left( \begin{array}{c c} 0 & 1 \\ 1 & 0 \end{array} \right) = \left( \begin{array}{c c} 1 & 0 \\ 0 & 1 \end{array} \right), \quad \left( \begin{array}{c c} 0 & 1 \\ 1 & 0 \end{array} \right) \left( \begin{array}{c c} a & 0 \\ 0 & b \end{array} \right) = \left( \begin{array}{c c} b & 0 \\ 0 & a \end{array} \right) \left( \begin{array}{c c} 0 & 1 \\ 1 & 0 \end{array} \right),
\]

this homomorphism extends to an algebra homomorphism \(\theta : \mathrm{F} \longrightarrow \mathbf{M}_2(\mathrm{K})\) defined by

\[
\theta \bigl ((a, b) + (c, d) j \bigr) = \left( \begin{array}{c c} a & c \\ d & b \end{array} \right)  .
\]

This homomorphism is bijective. When the characteristic of \(\mathbf{K}\) is different from 2, the algebra \(\mathbf{M}_2(\mathbf{K})\) is also isomorphic to the quaternion algebra of type (1,0,1) (Remark 1).

